% Part: intuitionistic-logic
% Chapter: tableaux
% Section: rules

\documentclass[../../../include/open-logic-section]{subfiles}

\begin{document}

\olfileid[pt]{int}{tab}{rul}

\olsection{Regras para a lógica intuicionista}

As regras para os conectivos $\land$ e $\lor$ são as mesmas que para
!!{tableau}s proposicionais sinalizados usuais, apenas com o acréscimo de prefixos.
Em cada caso, a regra aplicada a uma !!{formula} sinalizada
$\sFmla{S}{!A}[\sigma]$ produz novas !!{formula}s que também têm
o prefixo~$\sigma$. Isso deve ser intuitivamente claro: p.ex., se $!A
\land !B$ é verdadeira em um mundo nomeado por~$\sigma$, então $!A$ e $!B$
são verdadeiras em~$\sigma$ (e não em algum outro mundo). Reunimos as
regras para $\land$ e $\lor$ em \olref{tab:prop-rules}.

\begin{table}
  \[\def\arraystretch{3}\begin{array}{|c|c|}
    \hline
    \AxiomC{\sFmla{\True}{!A \land !B}[\sigma]}
    \RightLabel{\TRule{\True}{\land}}
    \UnaryInfC{\sFmla{\True}{!A}[\sigma]}
    \noLine
    \UnaryInfC{\sFmla{\True}{!B}[\sigma]}
    \DisplayProof
    &
    \AxiomC{\sFmla{\False}{!A \land !B}[\sigma]}
    \RightLabel{\TRule{\False}{\land}}
    \UnaryInfC{$\sFmla{\False}{!A}[\sigma] \quad \mid \quad
      \sFmla{\False}{!B}[\sigma]$}
    \DisplayProof
    \\[2ex]
    \hline
    \AxiomC{\sFmla{\True}{!A \lor !B}[\sigma]}
    \RightLabel{\TRule{\True}{\lor}}
    \UnaryInfC{$\sFmla{\True}{!A}[\sigma] \quad \mid \quad
      \sFmla{\True}{!B}[\sigma]$}
    \DisplayProof
    &
    \AxiomC{\sFmla{\False}{!A \lor !B}[\sigma]}
    \RightLabel{\TRule{\False}{\lor}}
    \UnaryInfC{\sFmla{\False}{!A}[\sigma]}
    \noLine
    \UnaryInfC{\sFmla{\False}{!B}[\sigma]}
    \DisplayProof
    \\[2ex]
    \hline
  \end{array}\]
  \caption{Regras prefixadas de !!{tableau} para $\land$ e $\lor$}
  \ollabel{tab:prop-rules}
\end{table}

A condição de fechamento é semelhante à dos !!{tableau}s usuais,
embora exijamos que não apenas as !!{formula}s, mas também os
prefixos coincidam. Na verdade, podemos ser um pouco mais permissivos: como,
nos modelos intuicionistas, as !!{formula}s permanecem verdadeiras uma vez que se tornam verdadeiras, é
impossível que $!A$ seja verdadeira em~$\sigma$, mas falsa em algum
prefixo acessível~$\sigma.{*}$. Portanto, um ramo é fechado se contém ambos
\[
\sFmla{\True}{!A}[\sigma] \quad\text{e}\quad \sFmla{\False}{!A}[\sigma.{*}]
\]
para algum prefixo $\sigma$ e alguma !!{formula}~$!A$. Observe que, se os sinais
estiverem invertidos, O.s.\@, se o ramo contiver
\[
\sFmla{\False}{!A}[\sigma] \quad\text{e}\quad \sFmla{\True}{!A}[\sigma.{*}]
\]
o ramo será fechado apenas se $*$ for a sequência vazia.

Além disso, um ramo é fechado se contém~$\sFmla{\True}{\bot}[\sigma]$.

As regras para estabelecer as premissas também são as mesmas que para
!!{tableau}s usuais, exceto que sempre usamos o
prefixo~$1$ para as premissas. (Não importa qual prefixo usamos, desde que seja
o mesmo para todas as premissas.) Assim, p.ex., dizemos que
\[
!B_1, \dots, !B_n \Proves !A
\]
se, e somente se, existe um tableau fechado para as premissas
\[
\sFmla{\True}{!B_1}[1], \dots, \sFmla{\True}{!B_n}[1],
\sFmla{\False}{!A}[1].
\]

Para o condicional~$\lif$, as regras diferem dos casos clássico e
modal. A regra $\TRule{\lif}{\True}$ estende um ramo
que contém $\sFmla{\True}{!A \lif !B}[\sigma]$ com
$\sFmla{\True}{!A}[\sigma.{*}]$ e $\sFmla{\False}{!B}[\sigma.{*}]$ em dois
ramos diferentes. Ela só pode ser aplicada a um prefixo~$\sigma.{*}$
que \emph{já} ocorra no ramo ao qual é aplicada.
Chamemos tal prefixo de ``usado'' no ramo. (Como $\sigma.{*}$
inclui o próprio $\sigma$, a regra sempre pode ser aplicada acrescentando as
fórmulas sinalizadas prefixadas $\sFmla{\True}{!A}[\sigma]$ e
$\sFmla{\False}{!B}[\sigma]$ em ramos separados.)

A regra $\TRule{\lif}{\False}$ estende um ramo que contém
$\sFmla{\False}{!A \lif !B}[\sigma]$ com
$\sFmla{\True}{!A}[\sigma.n]$ e $\sFmla{\False}{!B}[\sigma.n]$ no
mesmo ramo, sendo $\sigma.n$ um prefixo novo no ramo. 

As regras para $\lnot$ são definidas de modo análogo (usando a definição de
$\lnot !A$ como $!A \lif \lfalse$).

As regras são apresentadas em \olref{tab:rules-lif-lnot}.

\begin{table}
  \[\def\arraystretch{3}\begin{array}{|c|c|}
    \hline
    \AxiomC{\sFmla{\True}{\lnot !A}[\sigma]}
    \RightLabel{\TRule{\True}{\lnot}}
    \UnaryInfC{$\sFmla{\False}{!A}[\sigma.{*}]$}
    \DisplayProof
    &
    \AxiomC{\sFmla{\False}{\lnot !A}[\sigma]}
    \RightLabel{\TRule{\False}{\lnot}}
    \UnaryInfC{\sFmla{\True}{!A}[\sigma.n]}
    \DisplayProof
    \\[1ex]
    \text{$\sigma.{*}$ é usado} & \text{$\sigma.n$ é novo}\\
    \hline
    \AxiomC{\sFmla{\True}{!A \lif !B}[\sigma]}
    \RightLabel{\TRule{\True}{\lif}}
    \UnaryInfC{$\sFmla{\False}{!A}[\sigma.{*}] \quad \mid \quad
      \sFmla{\True}{!B}[\sigma.{*}]$}
    \DisplayProof
    &
    \AxiomC{\sFmla{\False}{!A \lif !B}[\sigma]}
    \RightLabel{\TRule{\False}{\lif}}
    \UnaryInfC{\sFmla{\True}{!A}[\sigma.n]}
    \noLine
    \UnaryInfC{\sFmla{\False}{!B}[\sigma.n]}
    \DisplayProof
    \\[1ex]
    \text{$\sigma.{*}$ é usado} & \text{$\sigma.n$ é novo}\\
    \hline
  \end{array}\]
  \caption{Regras prefixadas de !!{tableau} para $\lnot$ e $\lif$}
  \ollabel{tab:rules-lif-lnot}
\end{table}

\end{document}
